\documentclass[11pt,a4paper]{article}

\usepackage[utf8]{inputenc}
\usepackage[T1]{fontenc}
\usepackage{lmodern}
\usepackage{microtype}
\usepackage[margin=1in]{geometry}
\usepackage{hyperref}
\usepackage{graphicx}
\usepackage{booktabs}
\usepackage{amsmath}
\usepackage{amssymb}
\usepackage{xcolor}
\usepackage[backend=biber,style=numeric-comp,sorting=none]{biblatex}

\addbibresource{references.bib}

\hypersetup{
  colorlinks=true,
  linkcolor=black,
  citecolor=blue!50!black,
  urlcolor=blue!50!black,
  pdftitle={Verifiable Cross-Chain Risk Intelligence: OmniRisk, RayaChain, and Bittensor},
  pdfauthor={OmniRisk Research},
}

\title{Verifiable Cross-Chain Risk Intelligence:\\
       An Agentic Architecture over LayerZero V2 and Bittensor}
\author{OmniRisk Research \\ \texttt{research@omnirisk.io}}
\date{May 2026 \\ \textbf{Working paper, v0.3}}

\begin{document}
\maketitle

\begin{abstract}
Cross-chain bridges have become the single largest concentrated source of value loss in
decentralised finance: industry trackers attribute roughly USD~2B of theft in the first
eight months of 2022 alone to bridge attacks~\cite{chainalysis2022bridges}, and the most
recent peer-reviewed systematisation places cumulative losses at approximately
USD~3.1B~\cite{augusto2024sok}. At the same time, the risk telemetry that participants
need to navigate this environment remains \emph{siloed per chain} and
\emph{non-verifiable}: there is no widely deployed mechanism that lets a downstream user
confirm that a route- or token-level risk signal was produced honestly. We present
\textsc{OmniRisk}, a deployed agentic risk-intelligence layer that ingests cross-chain
flow data from \textsc{RayaChain}, an Omnichain Fungible Token (OFT) deployment on
LayerZero~V2~\cite{zarick2024layerzero,layerzero2024oft}, and emits structured route and
token risk signals. Our contribution is an architecture in which (i)~OmniRisk is the
intelligence layer, (ii)~RayaChain is the cross-chain data and execution layer, and
(iii)~Bittensor's Yuma Consensus~\cite{rao2021bittensor,steeves2022incentivizing,
opentensor2024yuma} acts as the decentralised verification layer for the intelligence
itself. We argue that anchoring the verification layer on the Subtensor EVM
(chain~ID~945)~\cite{opentensor2024evm} is the natural fit because it is the only
production EVM whose validator set is already economically aligned to scoring the
quality of AI outputs. The paper formalises the miner and validator protocols required
to make this work, situates the design against the recent cross-chain bridge security
literature~\cite{augusto2024sok,zhang2023sok,wu2025safeguarding}, and lays out the
evaluation framework we will use against historical RayaChain flows.
\end{abstract}

\section{Introduction}\label{sec:intro}

\textsc{OmniRisk} is a live cross-chain risk intelligence service, and \textsc{RayaChain}
is an actively deploying cross-chain transfer fabric. RayaChain's transfer lanes are
implemented as Omnichain Fungible Tokens (OFTs) on LayerZero~V2, the protocol family
that separates message verification (Decentralised Verifier Networks, DVNs) from message
execution (Executors) and exposes per-application security configuration as a
permissionless surface~\cite{zarick2024layerzero,layerzero2024oft}. As of submission,
RayaChain is operating across seven EVM and non-EVM ecosystems --- Ethereum, Arbitrum,
Base, BNB Smart Chain, Polygon, Optimism, and Solana --- and is scoping an eighth lane
into the Subtensor EVM (chain~ID~945)~\cite{opentensor2024evm}.

The motivation for the eighth lane is not capacity but \emph{verifiability}. Routing
value across heterogeneous chains exposes participants to a class of risks that survey
work has documented in detail: design flaws in bridge architectures map onto a stable
catalogue of attack surfaces~\cite{zhang2023sok}, attack-transaction analyses confirm
that the same patterns recur across dozens of incidents~\cite{wu2025safeguarding}, and
losses concentrate in the bridge layer to a degree that no other DeFi primitive
matches~\cite{chainalysis2022bridges,augusto2024sok}. Existing risk tools react to this
by emitting black-box risk scores: a downstream contract or wallet must trust the
operator's API, with no cryptographic or economic mechanism that constrains the
operator to be honest. This is the same primitive failure that gave rise to the
oracle-design literature in the early DeFi era, and we argue it now calls for a
verification layer that is itself decentralised.

\paragraph{Contributions.} This paper makes four contributions.
\begin{enumerate}
  \item \textbf{An architectural separation of concerns}: OmniRisk = intelligence,
        RayaChain = data and execution, Bittensor = verification (Section~\ref{sec:arch}).
  \item \textbf{A miner protocol specification} for thin OmniRisk-API wrappers that meet
        the determinism and structured-reasoning requirements of Yuma
        scoring~\cite{opentensor2024yuma} (Section~\ref{sec:miner}).
  \item \textbf{A validator-side scoring rule} that evaluates miner outputs against
        ground-truth realised cross-chain risk events, expressed in the Yuma weight-vector
        format (Section~\ref{sec:validator}).
  \item \textbf{An evaluation framework} grounded in replayed historical RayaChain
        flows, with an honest discussion of what we cannot yet prove
        (Section~\ref{sec:eval} and Section~\ref{sec:limitations}).
\end{enumerate}

This is a working paper. We are explicit throughout about which components are deployed
today, which are in active integration, and which are design proposals (see
Table~\ref{tab:status} in Section~\ref{sec:status}).

\section{Problem Definition}\label{sec:problem}

\paragraph{Unified risk visibility across chains is missing.} A token's risk profile is
not a property of any single chain. The same OFT can have safe liquidity on Ethereum
and dangerously thin liquidity on a destination rollup; a user routing through it will
realise the worse of the two. Empirical AMM-level evidence shows that liquidity is
materially mis-allocated across L1 and L2s relative to LP-return
equilibrium~\cite{gogol2024liquidity}, which means \emph{route choice} --- not just
token choice --- is a first-order risk decision.

\paragraph{Real-time route risk awareness is missing.} The bridge security literature
documents 49~attack incidents on cross-chain bridges between June~2021 and
September~2024~\cite{wu2025safeguarding}, and a security/privacy taxonomy spanning
57~interoperability solutions identifies recurring failure modes across them
all~\cite{augusto2024sok}. Yet at the moment a user submits a transfer, they typically
have no signal that the lane they are about to use shares its design with a recently
exploited one.

\paragraph{Verifiable intelligence is missing.} Existing route-risk APIs are
\emph{black-box}: a downstream contract must trust that the API operator scored honestly.
Recent work on optimistic and zero-knowledge ML on blockchain
(\cite{conway2024opml}, among others) and on multichain ML-based fraud
detection~\cite{palaiokrassas2023multichain} addresses parts of this gap, but no system
that we are aware of combines (a)~live cross-chain risk scoring with (b)~a decentralised
verification layer for the scores themselves.

\section{System Architecture}\label{sec:arch}

We split responsibility across three layers (Figure omitted; see Section~\ref{sec:eval}
for the empirical setup the architecture is designed to enable).

\subsection{OmniRisk --- Agentic Intelligence Layer}\label{sec:arch:omnirisk}

OmniRisk is a deployed system that observes cross-chain liquidity health, scores route
risk, and surfaces anomalous on-chain activity. We describe it as \emph{agentic} in the
narrow sense used in recent LLM systems work: a control loop that ingests chain state,
selects which downstream tools (DEX quote APIs, on-chain analytics endpoints,
historical incident databases) to consult, and emits a structured response with an
explicit confidence and a chain-of-evidence. The agentic framing is contrasted with
\emph{static} analytics --- precomputed dashboards, fixed risk-scoring formulas --- which
do not adapt to the question being asked.

\subsection{RayaChain --- Execution and Data Layer}\label{sec:arch:rayachain}

RayaChain is the OFT lane fabric. Each lane is implemented as either a standard OFT
(burn on source, mint on destination) or an OFT Adapter (lock on source, mint on
destination) where a legacy ERC-20 already exists, exactly as specified in the OFT
standard~\cite{layerzero2024oft}. RayaChain inherits LayerZero~V2's permissionless
security configuration: each lane chooses its own DVN set and execution path, which is
the surface OmniRisk monitors when scoring lane-level risk~\cite{zarick2024layerzero}.

The current deployment covers seven chains (Ethereum, Arbitrum, Base, BNB Smart Chain,
Polygon, Optimism, Solana). The proposed eighth lane is the Subtensor EVM
(chain~ID~945)~\cite{opentensor2024evm}, motivated in
Section~\ref{sec:arch:bittensor}.

\subsection{Bittensor --- Verification Layer}\label{sec:arch:bittensor}

Bittensor is a decentralised network in which \emph{miners} produce outputs for a
specific task and \emph{validators} score them; rewards are emitted via Yuma Consensus
on the Subtensor blockchain~\cite{rao2021bittensor,steeves2022incentivizing,
opentensor2024yuma}. Yuma takes each validator's stake-weighted weight vector over
miners, clips outliers against a stake-weighted threshold (consensus parameter
$\kappa = 0.5$), and emits a stake-aligned reward. The mechanism is provably
collusion-resistant up to coalitions controlling less than 50\% of stake-weighted
rank~\cite{steeves2022incentivizing}.

We use Bittensor as a \emph{verification} layer in the sense that miner risk scores are
not trusted on their face --- they are scored against validators' independent assessments,
and over time only miners whose scores correlate with the validator consensus survive
economically. This is the property we want from a risk oracle, and it is not provided by
any centralised risk API. The Subtensor EVM (chain~ID~945) gives us an Ethereum-compatible
execution environment in which to anchor commitments to score outputs without leaving
the Bittensor trust domain~\cite{opentensor2024evm}.

\section{OmniRisk--Bittensor Integration Model}\label{sec:integration}

The functional mapping is:
\begin{itemize}
  \item \textbf{OmniRisk} generates risk scores and anomaly signals (intelligence).
  \item \textbf{Bittensor validators} act as independent risk-verification nodes
        (verification).
  \item \textbf{RayaChain} contributes the cross-chain flow data and is also one of
        the execution venues whose lanes are being scored (data + execution).
\end{itemize}

A subnet specialised for cross-chain financial intelligence exposes a single task to
its miners: given a route specification (source chain, destination chain, OFT, size,
deadline), return a structured risk score. We do not require miners to all use the same
internal model --- the discipline comes from the validator scoring rule
(Section~\ref{sec:validator}), not from the miner implementation.

\section{Miner Protocol Specification}\label{sec:miner}

A miner is a thin wrapper over the OmniRisk API, exposed at the Bittensor subnet's
miner endpoint. For each query $q = (\textit{src}, \textit{dst}, \textit{oft},
\textit{size}, t)$ submitted by a validator, the miner returns a tuple
$(s, c, r) \in [0,1] \times [0,1] \times \mathcal{R}$ where:
\begin{itemize}
  \item $s$ is the route risk score, with $s = 0$ denoting nominal and $s = 1$ denoting
        an unrouteable lane;
  \item $c$ is the miner's self-reported confidence;
  \item $r$ is a structured reasoning record (which on-chain observations were used,
        which DEX quotes were consulted, which historical incidents were the closest
        analogues).
\end{itemize}

\paragraph{Determinism constraint.} Yuma scoring requires that two miners using the
same internal logic on the same query at the same block height return the same answer.
Otherwise validators cannot meaningfully aggregate them, and good miners are penalised
for irreducible randomness rather than for being wrong. We pin this by (i)~quantising
$s$ and $c$ to a fixed precision (we use 1/256), (ii)~anchoring the miner to a specific
chain head (a 32-byte block hash per chain consulted), and (iii)~specifying the
ordering and version of any external data sources in $r$. This is the same discipline
that makes optimistic-ML rollup systems
~\cite{conway2024opml} workable: outputs must be a deterministic function of the
declared inputs.

\section{Validator Mechanism}\label{sec:validator}

A validator queries a subset of miners on a stratified sample of routes, scores each
response against a ground-truth signal, and submits a Yuma weight vector
$\mathbf{w} \in [0,1]^{N}$ over the $N$ miners it observed. The on-chain Yuma rule
then aggregates $\mathbf{w}$ across all subnet validators using a stake-weighted
threshold and emits rewards~\cite{opentensor2024yuma}.

The validator scoring rule has three components:
\begin{enumerate}
  \item \textbf{Realised-event accuracy.} For routes that were taken between $t$ and
        $t + \Delta$, did the miner's score $s$ correlate with the realised loss
        (slippage, failed-execution rate, post-hoc identification as part of an attack
        cluster)?
  \item \textbf{Consistency.} If the same query is issued at $t$ and $t + \delta$ for
        small $\delta$ during which no chain state changed materially, do the responses
        match? Inconsistency indicates either non-determinism (penalty per
        Section~\ref{sec:miner}) or unstable internal state.
  \item \textbf{Confidence calibration.} Over a window of queries, does the miner's
        empirical Brier score on $(s, c)$ improve on a calibrated baseline? Miners that
        report $c = 1.0$ on every query are penalised relative to miners whose
        confidence tracks their realised accuracy.
\end{enumerate}

The key design property is that validators score the \emph{quality} of intelligence, not
just its presence. Miners that route every score through the same underlying OmniRisk
API but apply different downstream models can still be differentiated, because the
realised-event channel breaks ties.

\section{Cross-Chain Risk Model}\label{sec:risk}

We formalise three risk axes that emerge \emph{between} chains rather than within them:
\begin{description}
  \item[Liquidity fragmentation.] The same token has independent depth on each chain;
        the worst-of distribution dominates a multi-hop route. The empirical
        misallocation documented by~\cite{gogol2024liquidity} for WETH--USDC pools is
        the simplest instance of this; the OFT lane fabric of LayerZero~V2 generalises
        the problem to arbitrary fungible tokens~\cite{layerzero2024oft}.
  \item[Bridge / route dependency risk.] A given OFT lane inherits the security
        configuration of its DVN set and Executor~\cite{zarick2024layerzero}. The
        bridge-security literature catalogues how individual configuration choices map
        onto exploited vulnerability classes
        ~\cite{zhang2023sok,augusto2024sok,wu2025safeguarding}.
  \item[Anomaly propagation.] An exploit on one chain can spread to others via shared
        bridges, shared liquidity, and shared dependencies. Multichain fraud-detection
        datasets~\cite{palaiokrassas2023multichain} provide the empirical basis for
        modelling this, but the cross-chain dimension remains under-instrumented in
        published work.
\end{description}
The unifying observation is that cross-chain risk is \emph{a property of routes, not of
tokens or chains in isolation}, and a verification layer that scores claims about routes
is therefore the natural primitive.

\section{Subtensor EVM (Chain~ID~945) Integration}\label{sec:subtensor}

The Subtensor EVM is an Ethereum-compatible execution environment that runs inside the
Subtensor Substrate runtime, with TAO as the gas asset~\cite{opentensor2024evm}. It is
not just another EVM lane for RayaChain; it is the execution environment whose
validator set is already economically aligned to scoring the quality of off-chain
intelligence, by virtue of the Yuma mechanism~\cite{opentensor2024yuma,
steeves2022incentivizing}. We use it for two distinct purposes:
\begin{enumerate}
  \item \textbf{Anchoring of risk-score commitments.} Each batch of miner outputs is
        committed as a Merkle root on Subtensor EVM, so a downstream contract on any
        of the seven existing RayaChain lanes can fetch a proof that a given
        $(s, c, r)$ tuple was produced and made available at a specific block height,
        prior to the route being taken.
  \item \textbf{Anchoring of validator-consensus outputs.} The aggregated Yuma reward
        vector for the cross-chain risk subnet is itself published on Subtensor EVM,
        which gives a downstream auditor a single canonical record of which miners
        were trusted at any given epoch.
\end{enumerate}
This aligns the economic incentives of the AI layer (miners and validators), the
liquidity layer (RayaChain LPs), and the verification layer (TAO holders staking on
the subnet).

\section{Data and Evaluation Framework}\label{sec:eval}

The evaluation we are building rests on three datasets:
\begin{enumerate}
  \item \textbf{RayaChain flow logs.} OFT send/receive events across the seven live
        lanes, with per-lane metadata (DVN set, Executor, observed latencies,
        message-failure counts).
  \item \textbf{Realised-event labels.} A union of (a)~slippage and execution-failure
        outcomes derived from the flow logs, and (b)~external incident labels from
        peer-reviewed and grey-literature corpora; the BridgeGuard
        dataset~\cite{wu2025safeguarding} (203 attack transactions across
        49~incidents) is one of the curated bases.
  \item \textbf{Replay routes.} For each historical route observed in (1), we replay
        it to the OmniRisk-backed miners under the configuration at the corresponding
        block height (Section~\ref{sec:miner}), and score against (2).
\end{enumerate}
The headline metrics are realised-event AUC, early-detection lead time (how many
blocks before an event the score crossed threshold), and Yuma reward retention (a
miner's reward share over a window normalised by its scored share, as a measure of
alignment with the validator consensus).

\section{Economic Model}\label{sec:economics}

The economic flows are:
\begin{itemize}
  \item Miners are rewarded in TAO via Yuma emissions for risk outputs that the
        validator consensus accepts~\cite{opentensor2024yuma,steeves2022incentivizing}.
  \item Validators are rewarded for assigning weights that align with the consensus,
        which is the natural Yuma incentive to score accurately.
  \item OmniRisk is the infrastructure provider: it operates the underlying analytics
        platform that miners wrap, and it monetises this through its existing API
        surface to non-Bittensor consumers.
\end{itemize}
We avoid claims about token-economic equilibria here. The Bittensor critical-analysis
literature~\cite{lui2025bittensor} documents real concerns about Bittensor's
centralisation properties under dynamic TAO, and any quantitative economic claim about
this subnet should be revisited against that work in a follow-up.

\section{Security Model}\label{sec:security}

The threat surface is dominated by three classes:
\begin{description}
  \item[Adversarial miners.] A miner can attempt to bias the risk score it emits, for
        example to suppress a warning on a lane it benefits from. The Yuma scoring rule
        bounds the influence of such a miner to its share of stake-weighted
        rank~\cite{steeves2022incentivizing}; the determinism and consistency
        components of the validator rule (Section~\ref{sec:validator}) further limit
        steady-state bias.
  \item[Manipulated liquidity signals.] An attacker can briefly distort the on-chain
        signals that miners ingest --- e.g.\ wash-trading a thin pool to mask an exit.
        Our mitigation is multi-source data validation: the validator's
        realised-event channel (Section~\ref{sec:validator}) is by construction
        retrospective and ignores manipulated quote-time signals.
  \item[Cross-chain spoofing.] A malicious bridge can claim a state on one chain that
        is not consistent with the source. The OFT design~\cite{layerzero2024oft}
        and the V2 DVN model~\cite{zarick2024layerzero} reduce, but do not eliminate,
        this risk; the bridge-security literature documents the residual surface in
        depth~\cite{zhang2023sok,augusto2024sok,wu2025safeguarding}.
\end{description}

\section{Implementation Status}\label{sec:status}

We separate what is deployed from what is in design (Table~\ref{tab:status}).

\begin{table}[h]
\centering
\begin{tabular}{lll}
\toprule
\textbf{Component} & \textbf{Status as of 2026-05} & \textbf{Reference} \\
\midrule
OmniRisk intelligence layer        & Live in production       & this paper \\
RayaChain lanes (7 chains)         & Live / actively deploying & \cite{zarick2024layerzero,layerzero2024oft} \\
Subtensor EVM lane (chain~945)     & Planned (8th lane)        & \cite{opentensor2024evm} \\
Cross-chain risk subnet (Bittensor)& Design proposal           & this paper \\
Validator scoring rule             & Design proposal           & \S\ref{sec:validator} \\
Score-commitment anchoring on TAO  & Design proposal           & \S\ref{sec:subtensor} \\
\bottomrule
\end{tabular}
\caption{Implementation status. Components marked \emph{design proposal} are
specified in this paper but not yet deployed, and quantitative claims about them are
deferred to follow-up work.}
\label{tab:status}
\end{table}

\section{Roadmap}\label{sec:roadmap}

The work plan implied by Table~\ref{tab:status} is, in order:
(i)~finalise the miner API and the determinism contract (Section~\ref{sec:miner});
(ii)~stand up the cross-chain risk subnet on Subtensor and integrate the validator
scoring rule (Section~\ref{sec:validator});
(iii)~deploy the score-commitment contracts on Subtensor EVM (Section~\ref{sec:subtensor});
(iv)~ship the eighth RayaChain lane into Subtensor EVM;
(v)~publish the evaluation results from Section~\ref{sec:eval} as a follow-up paper.

\section{Comparative Advantage}\label{sec:compare}

Existing cross-chain risk tooling is some combination of single-chain, static, or
black-box. The architecture we propose is differentiated on three axes simultaneously:
\begin{itemize}
  \item \textbf{Agentic, not static.} The intelligence layer adapts its tool use to the
        question, rather than precomputing a fixed risk surface.
  \item \textbf{Cross-chain-native, not single-chain.} Routes, not tokens, are the
        primitive being scored, which is the level at which fragmentation and
        bridge-dependency risk actually emerge~\cite{gogol2024liquidity,
        zhang2023sok}.
  \item \textbf{Verifiable, not black-box.} The Bittensor verification layer makes
        miner outputs economically expensive to bias, and the on-chain commitments on
        Subtensor EVM (Section~\ref{sec:subtensor}) make their history auditable from
        any of the seven other lanes.
\end{itemize}

\section{Limitations}\label{sec:limitations}

We explicitly flag the following limitations.
\begin{itemize}
  \item \textbf{Most quantitative claims are deferred.} The score-commitment
        contracts, the Bittensor subnet, and the evaluation framework
        (Section~\ref{sec:eval}) are design proposals; this paper does not yet report
        AUC numbers or reward-retention curves.
  \item \textbf{Bittensor centralisation risk.} Independent empirical work has raised
        concerns about Bittensor's actual decentralisation under the dynamic-TAO
        regime~\cite{lui2025bittensor}. We treat the Yuma collusion-resistance bound
        of \cite{steeves2022incentivizing} as a \emph{design property of the algorithm}
        rather than as a guarantee about the deployed network's stake distribution.
  \item \textbf{Withdrawn primary source.} The original BitTensor preprint
        (arXiv:2003.03917) was withdrawn by its authors as obsolete on 2021-11-10. Any
        prior whitepaper that cited it should be reread against the canonical
        \cite{rao2021bittensor} or the workshop
        formalisation~\cite{steeves2022incentivizing}.
  \item \textbf{Determinism in agentic LLM components.} Section~\ref{sec:miner}'s
        determinism contract is straightforward for classical analytics but is more
        delicate when the miner uses an LLM under the hood; the
        consistency component of the validator rule
        (Section~\ref{sec:validator}) is what catches the residual non-determinism in
        practice.
\end{itemize}

\section{Conclusion}\label{sec:conclusion}

Cross-chain risk has settled into a regime where the dominant losses come from a small
number of well-understood failure modes in the bridge layer, but where the risk
\emph{intelligence} that participants need is itself untrusted. We have presented an
architecture that splits intelligence (OmniRisk), data and execution (RayaChain on
LayerZero~V2), and verification (Bittensor's Yuma Consensus, anchored on Subtensor EVM
chain~ID~945) into three layers with clean economic boundaries. The result, when fully
deployed, is the first production cross-chain risk primitive whose outputs are
verifiable in the same sense that the underlying transfers are: by an independent,
economically aligned, decentralised process. The intermediate deliverables on the
roadmap (Section~\ref{sec:roadmap}) are testable in isolation, and we will report on
each in follow-up work.

\section*{Acknowledgements}
We thank the OmniRisk engineering team for the production telemetry that grounds
Section~\ref{sec:eval}, and the Opentensor Foundation maintainers for the Subtensor
EVM testnet documentation that made the Section~\ref{sec:subtensor} integration plan
concrete. Any errors are the authors'.

\section*{Reproducibility Note}
The bibliography for this paper is the master \texttt{references.bib} of the
\texttt{omnirisk-papers} repository. Every entry was reverified against its primary
source on 2026-05-03 prior to inclusion. The retracted preprint arXiv:2003.03917
is intentionally excluded; see the comment block in \texttt{references.bib} and
Section~\ref{sec:limitations} for the rationale.

\printbibliography

\end{document}
